% Figure for Electromagnetism §A6.3 (Example: a two-loop circuit): assumed currents and the two loops used.
\begin{circuitikz}[>={Stealth[length=2.0mm]}, font=\small]
  \draw (0,3) to[R, l={$R_1 = 3\ \Omega$}] (3,3) to[R, l={$R_3 = 7\ \Omega$}] (6,3) -- (6,0)
        to[battery1, l={$V_2 = 29\ \mathrm V$}] (3,0) -- (0,0);
  \draw (0,3) to[battery1, l_={$V_1 = 24\ \mathrm V$}] (0,0);
  \draw (3,3) to[R, l_={$R_2 = 3\ \Omega$}] (3,0);
  \fill (3,3) circle (1.6pt) (3,0) circle (1.6pt);
  % point labels
  \node[font=\scriptsize, anchor=south east] at (0,3) {$a$};
  \node[font=\scriptsize, anchor=south] at (3,3.05) {$b$};
  \node[font=\scriptsize, anchor=south west] at (6,3) {$c$};
  \node[font=\scriptsize, anchor=north west] at (6,0) {$d$};
  \node[font=\scriptsize, anchor=north] at (3,-0.05) {$e$};
  \node[font=\scriptsize, anchor=north east] at (0,0) {$f$};
  \node[font=\scriptsize] at (-0.3,1.95) {$+$};
  \node[font=\scriptsize] at (4.95,0.35) {$+$};
  % assumed currents
  \draw[->, thick, figR] (1.9,2.6) -- (2.6,2.6) node[midway, below, font=\scriptsize] {$I_1$};
  \draw[->, thick, figR] (3.35,2.3) -- (3.35,1.6) node[midway, right, font=\scriptsize] {$I_2$};
  \draw[->, thick, figR] (3.4,2.6) -- (4.1,2.6) node[midway, below, font=\scriptsize] {$I_3$};
  % loops (clockwise)
  \draw[->, figB, thin] (0.85,1.3) arc[start angle=200, end angle=-120, radius=0.42];
  \draw[->, figB, thin] (4.75,1.3) arc[start angle=200, end angle=-120, radius=0.45];
\end{circuitikz}
